% Generated draft: requires mathematical and source review.
\documentclass[11pt]{article}
\usepackage[margin=25mm]{geometry}
\usepackage{fontspec}
\setmainfont{Latin Modern Roman}
\usepackage{amsmath,amssymb,mathtools,unicode-math}
\setmathfont{Latin Modern Math}
\usepackage{enumitem,needspace,xurl,hyperref}
\setlength{\parindent}{0pt}
\setlength{\parskip}{5pt}
\setlength{\emergencystretch}{3em}
\newcommand{\N}{\mathbb{N}}
\newcommand{\Z}{\mathbb{Z}}
\newcommand{\Q}{\mathbb{Q}}
\newcommand{\R}{\mathbb{R}}
\newcommand{\C}{\mathbb{C}}
\newcommand{\F}{\mathbb{F}}
\newcommand{\A}{\mathbb{A}}
\newcommand{\PP}{\mathbb{P}}
\newcommand{\E}{\mathbb{E}}
\newcommand{\eps}{\varepsilon}
\newcommand{\dd}{\,\mathrm{d}}
\begin{document}
\section*{UnsolvedMath ID: 11100047}
\textbf{Source problem code:} catalogue entry
\subsection{Definitions and mathematical statement}
Let \(X\) be a scheme. The source asks to understand and construct a model category on complexes of sheaves, tentatively quasi-coherent sheaves, and expects compatibility with tensor product. The category must be chosen: \(\operatorname{Ch}(\operatorname{Mod}\mathcal O_X)\) and \(\operatorname{Ch}(\operatorname{QCoh}(X))\) are different inputs, whose hypotheses and derived-category comparisons must be stated.

A precise proposed model structure should give weak equivalences, cofibrations and fibrations satisfying the model-category axioms, with weak equivalences the quasi-isomorphisms
\[f:C\to D,\qquad H_j(f)\text{ is an isomorphism for every }j,\]
and identify its homotopy category with the corresponding derived category. The bounded/unbounded convention is not stated in the source. For symmetric monoidality under the usual total tensor product of complexes, it must satisfy the pushout-product axiom: for cofibrations \(i:A\to B\), \(j:C\to D\),
\[(B\otimes C)\amalg_{A\otimes C}(A\otimes D)\to B\otimes D\]
is a cofibration, trivial when either input is trivial, together with the unit axiom. This ensures that the derived tensor product has the required model-theoretic compatibility; merely having a model structure is insufficient.

Hovey frames this as an example of his broader motivation to find model categories underlying applications of homology, not as a theorem that every unspecified homology construction has one. He says a sheaf model probably already exists and contrasts the scarcity of projectives with the injective model's general failure to be symmetric monoidal. Those are his expressed uncertainties and explanation of what he wishes to learn, not assumptions that all sheaf categories have enough projectives.

The supplied summary reports a solution via symmetric monoidal flat model structures, attributing it to Hovey (2001). Its blanket mention of projective structures in general Grothendieck/sheaf categories and its truncated monoidality description need hypotheses: injective-type existence, flat/cotorsion constructions and enough-projective arguments cannot be interchanged. The precise scheme, flat-generation and tensor hypotheses of a claimed construction must be matched to the cited result, rather than asserting one unrestricted model for every category listed.
\subsection{Short English statement}
Construct and explain a quasi-isomorphism model structure on the selected sheaf-complex category that is also symmetric monoidal. The record reports flat-model solutions, while the scheme/category hypotheses and blanket projective claims need checking.
\subsection{Sources}
\begin{itemize}
\item Local UnsolvedMath registry, permanent ID \(11100047\), exact source code \texttt{catalogue entry}.
\item Mark Hovey, \emph{Hovey's Algebraic Topology Problem List}, Model categories, item 1, surviving mirror: \url{https://www-users.cse.umn.edu/~tlawson/hovey/}. Registry access 2026-07-29; html-ordered-list extraction.
\item Hovey, \emph{Model category structures on chain complexes of sheaves} (2001), as supplied in the summary; exact publication and theorem/scheme hypotheses are not reproduced.
\end{itemize}
\end{document}
